\subsection{The KAN-ODE model}
\label{sec:architecture}

Every KAN-ODE we train has the same architecture: a two-layer network with
widths $[2,10,2]$, each edge carrying the expansion of
Equation~\ref{eq:kan-edge} over a grid of $G=5$ basis functions. It is
wrapped inside an explicit Runge-Kutta integrator that advances the state
and hands the resulting trajectory to the loss.
Figure~\ref{fig:kan-ode-architecture} shows the full training loop and marks
the three components we vary.

\begin{figure*}[t]
\centering
\includegraphics[width=\textwidth]{figures/tikz/kan_ode_system.pdf}
\caption{The KAN-ODE training loop as implemented. The integrator calls the
network once per Runge-Kutta stage, so one training epoch on $[0,3.5]$
evaluates $\mathbf{f}_\theta$ $70s$ times ($35$ sampling intervals, two
sub-steps each, $s$ stages per step). The loss gradient comes from
differentiating through every one of those evaluations. Numbered badges mark
the three components we vary: the edge basis
(Sections~\ref{sec:basis-results} and~\ref{sec:track-c}), the integrator
(Sections~\ref{sec:solver-results}, \ref{sec:track-b} and~\ref{sec:track-d})
and the gradient method (Section~\ref{sec:track-a}).}
\label{fig:kan-ode-architecture}
\end{figure*}

\subsubsection{One KAN layer as two matrix products}

The network is two KDense layers. For a batch
$\mathbf{X}\in\mathbb{R}^{B\times d_{\text{in}}}$, a layer evaluates every
edge function of Equation~\ref{eq:kan-edge} at once:
\begin{equation}
\Psi_{b,i,g} = B_g\big(\tanh X_{b,i}\big), \qquad
\mathbf{Y} = \underbrace{\Psi_{\text{flat}}\,\mathbf{C}^{\top}}_{\text{spline branch}}
           + \underbrace{\mathrm{silu}(\mathbf{X})\,\mathbf{W}^{\top}}_{\text{base branch}},
\label{eq:kdense}
\end{equation}
with $\mathbf{C}\in\mathbb{R}^{d_{\text{out}}\times d_{\text{in}}G}$ and
$\mathbf{W}\in\mathbb{R}^{d_{\text{out}}\times d_{\text{in}}}$. Here
$\Psi_{\text{flat}}\in\mathbb{R}^{B\times d_{\text{in}}G}$ is $\Psi$ with
its last two axes merged. Entry $C_{j,(i,g)}$ is the coefficient $c_g$ of
edge $i\to j$ and $W_{j,i}$ is that edge's residual weight $w_b$, so
$Y_{b,j}=\sum_i\phi_{i,j}(X_{b,i})$ exactly. Figure~\ref{fig:kdense} traces
the tensor shapes. The forward pass is short enough to show in full:

\begin{pycode}[Forward pass of one KDense layer]
x_flat = x.view(-1, self.in_features)          # [B, d_in]
x_norm = self.normalizer(x_flat)               # tanh -> [-1, 1]
basis_vals = self.basis_func(x_norm, self.grid, self.denominator)
basis_flat = basis_vals.reshape(basis_vals.size(0),
                                self.in_features * self.grid_len)
spline_out = F.linear(basis_flat, self.C)      # Psi_flat @ C^T
if self.use_base_act and self.W is not None:
    base_val = self.base_act(x_flat)           # silu(x)
    base_out = F.linear(base_val, self.W)      # silu(x) @ W^T
    out = spline_out + base_out
else:
    out = spline_out
return out.view(*orig_shape[:-1], self.out_features)
\end{pycode}

\begin{figure*}[t]
\centering
\includegraphics[width=0.9\textwidth]{figures/tikz/kdense_layer.pdf}
\caption{Tensor flow through one KDense layer
(Equation~\ref{eq:kdense}). Only the basis-expansion step differs between
the seven basis families of Section~\ref{sec:bases}; every basis function
has the same signature,
$(\mathbf{U},\mathbf{z},h)\mapsto\Psi\in\mathbb{R}^{B\times d_{\text{in}}\times G}$,
so swapping the basis changes one tensor and nothing else.}
\label{fig:kdense}
\end{figure*}

\paragraph{Parameter count.} Each layer holds
$d_{\text{out}}d_{\text{in}}G$ spline coefficients,
$d_{\text{out}}d_{\text{in}}$ residual weights and no bias. For $[2,10,2]$
with $G=5$, layer one ($2\to10$) has $10\cdot2\cdot5+10\cdot2=120$
parameters and layer two ($10\to2$) has $2\cdot10\cdot5+2\cdot10=120$, for
the $240$ total the base paper reports. The comparison MLPs of
Section~\ref{sec:kan-vs-mlp} are sized to match: the paper's $[2,50,2]$
network has $2\cdot50+50+50\cdot2+2=252$ parameters, and the deeper
$[2,14,8,8,2]$ network we use as the working baseline also has
$(28+14)+(112+8)+(64+8)+(16+2)=252$.

\paragraph{Why the \texttt{efficient-kan} formulation.} The reference
implementation, \texttt{pykan} \cite{liu2024kan}, represents every edge with
its own spline object and recomputes the grid during training. Inside a
Neural ODE that cost is multiplied by the integrator: one epoch of Tsit5 on
the training window calls $\mathbf{f}_\theta$ $420$ times, and reverse-mode
differentiation keeps every one of those calls in the autograd graph. Our
layer follows the \texttt{efficient-kan} reformulation
\cite{blealtan2024efficientkan} instead: the knots are a fixed buffer, the
basis is evaluated for all edges with one broadcast, and the contraction is
one dense matrix product per branch (Equation~\ref{eq:kdense}). Each stage
therefore adds two matrix multiplications and a handful of elementwise
operations to the graph, however many edges there are. The same structure
makes every basis we substitute into Equation~\ref{eq:kan-edge} (RBF,
B-spline, Chebyshev, Lagrange, Newton, IQF, RSWAF) a swappable tensor rather
than a different object type.

\paragraph{What we implemented ourselves.} We wrote the KAN layer, the seven
basis functions, the six fixed-step Runge-Kutta integrators, the
invariant-enforcing field wrappers used for SIR and the training loop
directly in PyTorch \cite{paszke2019pytorch}, rather than importing them, so
that every basis and solver could be instrumented in the same way. The
integrators apply the tableaux of Section~\ref{sec:integrators} verbatim.
Tsit5 and DOPRI5 run at a fixed step with their fifth-order weights only,
whereas the base paper uses Tsit5's embedded error estimate for adaptive
step control. The base paper compiles its integration loop in Julia's SciML
ecosystem; we run the same loop uncompiled in PyTorch because it was simpler
for us to work with. Our wall-clock figures are therefore comparable across
our own runs but not against the paper's reported times
(Section~\ref{sec:budget} quantifies what this costs).

\subsection{Runge-Kutta integrators}
\label{sec:integrators}

Every integrator we use is an explicit Runge-Kutta method: one step of size
$h$ evaluates $s$ stage slopes and combines them,
\begin{equation}
\mathbf{k}_i=\mathbf{f}_\theta\Big(\mathbf{u}_n+h\sum_{j=1}^{i-1}a_{ij}\mathbf{k}_j\Big),
\qquad
\mathbf{u}_{n+1}=\mathbf{u}_n+h\sum_{i=1}^{s}b_i\mathbf{k}_i,
\label{eq:rk-general}
\end{equation}
for $i=1,\dots,s$, so a method is fully specified by its Butcher tableau
$(A,\mathbf{b},\mathbf{c})$ with $c_i=\sum_j a_{ij}$. Writing
$\mathbf{u}_n$ for the state at step $n$, the four low-order methods are
\begin{align}
\text{Euler ($p=1$):}\ & \mathbf{u}_{n+1} = \mathbf{u}_n + h\, \mathbf{f}_\theta(\mathbf{u}_n) \label{eq:euler}\\
\text{Midpoint ($p=2$):}\ & \mathbf{u}_{n+1} = \mathbf{u}_n + h\,\mathbf{f}_\theta\!\left(\mathbf{u}_n + \tfrac{h}{2}\mathbf{k}_1\right) \label{eq:midpoint}\\
\text{Heun ($p=2$):}\ & \mathbf{u}_{n+1} = \mathbf{u}_n + \tfrac{h}{2}(\mathbf{k}_1+\mathbf{k}_2),\ \ \mathbf{k}_2 = \mathbf{f}_\theta(\mathbf{u}_n+h\,\mathbf{k}_1) \label{eq:heun}\\
\text{RK4 ($p=4$):}\ & \mathbf{u}_{n+1} = \mathbf{u}_n + \tfrac{h}{6}\left(\mathbf{k}_1+2\mathbf{k}_2+2\mathbf{k}_3+\mathbf{k}_4\right) \label{eq:rk4}
\end{align}
with $\mathbf{k}_1=\mathbf{f}_\theta(\mathbf{u}_n)$ throughout and the
classical RK4 stages
$\mathbf{k}_2=\mathbf{f}_\theta(\mathbf{u}_n+\tfrac{h}{2}\mathbf{k}_1)$,
$\mathbf{k}_3=\mathbf{f}_\theta(\mathbf{u}_n+\tfrac{h}{2}\mathbf{k}_2)$,
$\mathbf{k}_4=\mathbf{f}_\theta(\mathbf{u}_n+h\,\mathbf{k}_3)$. DOPRI5 and
Tsit5 are both seven-stage 5(4) embedded pairs whose seventh stage exists
only to supply the fourth-order error estimate (the ``first same as last''
stage); its fifth-order weight is $b_7=0$. At a fixed step, as here, the
seventh evaluation contributes nothing to the update, so our implementation
skips it and both methods cost \emph{six} evaluations per step. Their
coefficients satisfy all $17$ order conditions up to order five, and
Tsitouras \cite{tsitouras2011tsit5} additionally minimised the size of the
sixth-order error terms. Table~\ref{tab:tsit5} lists the Tsit5 coefficients
exactly as implemented, and Figure~\ref{fig:rk-stages} shows where each
method samples the field within one step. Every extra stage costs one more
evaluation of $\mathbf{f}_\theta$, which is why we measure cost in function
evaluations per epoch rather than by solver order alone.

\begin{table}[t]
\centering
\scriptsize
\setlength{\tabcolsep}{2.6pt}
\renewcommand{\arraystretch}{1.15}
\caption{Fixed-step Tsit5 tableau as implemented (first six stages, rounded
to six significant figures; the code carries 16). Row $i$ lists $c_i$ and
$a_{i1},\dots,a_{i,i-1}$; the last two rows list the solution weights $b_i$.}
\label{tab:tsit5}
\begin{tabular}{@{}r|rrrrr@{}}
$0$ & & & & & \\
$0.161$ & $0.161$ & & & & \\
$0.327$ & $-0.00848066$ & $0.335481$ & & & \\
$0.9$ & $2.89715$ & $-6.35945$ & $4.36230$ & & \\
$0.980026$ & $5.32586$ & $-11.7489$ & $7.49554$ & $-0.0924951$ & \\
$1$ & $5.86146$ & $-12.9210$ & $8.15937$ & $-0.0715850$ & $-0.0282691$\\
\hline
$\mathbf{b}^\top$ & $0.0964608$ & $0.01$ & $0.479890$ & $1.37901$ & $-3.29007$\\
 & $2.32471$ & & & & \\
\end{tabular}
\end{table}

\begin{figure}[t]
\centering
\includegraphics[width=\columnwidth]{figures/tikz/rk_stages.pdf}
\caption{Stage nodes $t_n+c_ih$ and solution weights $b_i$ of the six
implemented methods, drawn from the coefficients in the code. Heun and
Midpoint use the same number of evaluations but sample at opposite places:
Heun averages the two ends of the step, Midpoint uses only a slope at its
centre. Tsit5 reaches order five with large weights of both signs
($b_5=-3.29$), a cancellation that is harmless in exact arithmetic and buys
its small error constant.}
\label{fig:rk-stages}
\end{figure}

\subsection{Basis functions}
\label{sec:bases}

Writing $z_g$ for the $g$-th of $G$ evenly spaced grid centres with spacing
$h=2/(G-1)$, and $u=\tanh(x)\in[-1,1]$ for the normalized edge input, the
seven bases we compare are, exactly as implemented,
\begin{gather}
B_g^{\text{RBF}}(u)=\exp\!\left(-\tfrac12\left(\tfrac{u-z_g}{h}\right)^2\right),\quad
B_g^{\text{IQF}}(u)=\frac{1}{1+\left(\frac{u-z_g}{h}\right)^2}, \nonumber\\
B_g^{\text{RSWAF}}(u)=1-\tanh^2\!\left(\tfrac{u-z_g}{h}\right),
\label{eq:basis-rbf-spline}
\end{gather}
\begin{gather}
B_g^{\text{Cheb}}(u)=T_{g-1}(u)=\cos\big((g-1)\arccos u\big), \nonumber\\
B_g^{\text{Lag}}(u)=\prod_{j\neq g}\frac{u-z_j}{z_g-z_j}, \qquad
B_g^{\text{Newt}}(u)=\prod_{j<g}(u-z_j),
\label{eq:basis-poly}
\end{gather}
plus the cubic B-spline $B_g^{\text{spline}}$, built by the Cox--de Boor
recursion from degree-0 indicator functions by repeated linear blending,
\begin{gather}
B_{g,0}(u)=\mathbf{1}[t_g\le u<t_{g+1}], \nonumber\\
B_{g,p}(u)=\frac{u-t_g}{t_{g+p}-t_g}B_{g,p-1}(u)+\frac{t_{g+p+1}-u}{t_{g+p+1}-t_{g+1}}B_{g+1,p-1}(u),
\label{eq:cox-de-boor}
\end{gather}
on the clamped knot vector $\mathbf{t}=(-1,-1,-1,-1,0,1,1,1,1)$ for $G=5$,
$k=3$. The clamping makes the five splines a partition of unity with
$B_1(-1)=B_5(1)=1$, and it is what gives the spline compact support:
$B_{g,k}(u)=0$ outside the $k+2$ knots the recursion references. Gaussian
RBF, IQF and RSWAF are local but never exactly zero. Chebyshev, Lagrange and
Newton have no locality at all: every coefficient $c_g$ influences $\phi(x)$
at every $x$. The residual term $w_b\,\mathrm{silu}(x)$ in
Equation~\ref{eq:kan-edge} is the same for every basis, since it is a fixed
part of the edge parameterization rather than one of the substituted terms.

\subsection{Computing gradients: direct autograd and the adjoint method}
\label{sec:gradients}

There are two ways to compute gradients through a Neural ODE. \textbf{Direct
autograd} treats the numerical integrator as an ordinary computational
graph: each of the solver's $N_t$ forward steps is recorded, and the reverse
pass walks back through every recorded step, so peak memory grows linearly
with the number of steps. \textbf{Adjoint sensitivity}
\cite{chen2018neuralode} instead treats the forward solve as a black box and
derives a second, augmented ODE that produces the same gradient without
storing the intermediate states. Writing $\mathbf{a}(t) = \partial
L/\partial\mathbf{u}(t)$ for the adjoint state, it satisfies its own
backward-time differential equation,
\begin{equation}
\frac{d\mathbf{a}}{dt} = -\mathbf{a}^\top\frac{\partial \mathbf{f}_\theta}{\partial\mathbf{u}},
\qquad
\frac{dL}{d\theta} = \int_{t_1}^{t_0}\mathbf{a}^\top\frac{\partial\mathbf{f}_\theta}{\partial\theta}\,dt,
\label{eq:adjoint}
\end{equation}
solved backward from $t_1$ to $t_0$ alongside the original state
$\mathbf{u}(t)$, which is recomputed rather than stored. Memory then depends
only on the size of the state and parameter vectors, not on how many steps
the forward solve took: $O(1)$ in trajectory length, against direct
autograd's $O(N_t)$ (Figure~\ref{fig:adjoint-schematic}). All our training
runs use direct autograd; Section~\ref{sec:track-a} measures whether that
choice is justified.

\begin{figure}[t]
\centering
\resizebox{\columnwidth}{!}{%
\begin{tikzpicture}[
  font=\small,
  statend/.style={draw=ink, fill=accentsoft, minimum width=9mm, minimum height=7mm, rounded corners=2pt},
  adjnd/.style={draw=ink!60, fill=codebg, minimum width=11mm, minimum height=7mm, rounded corners=2pt},
  arr/.style={-{Stealth[length=2.2mm]}, draw=ink!70, thick},
  panellbl/.style={font=\small\sffamily\bfseries\color{ink}}
]
  \node[statend] (u0) at (0,0) {$u_0$};
  \foreach \i/\x in {1/1.5,2/3.0,3/4.5,4/6.0} {
    \node[statend] (u\i) at (\x,0) {$u_\i$};
  }
  \draw[arr] (u0) -- (u1); \draw[arr] (u1) -- (u2);
  \draw[arr] (u2) -- (u3); \draw[arr] (u3) -- (u4);
  \draw[{Stealth[length=2.2mm]}-, dashed, accent, thick]
    (u0.south) .. controls +(0,-0.9) and +(0,-0.9) .. (u4.south)
    node[midway, below, font=\scriptsize\sffamily\color{accent}] {backward pass walks every stored $u_i$};
  \node[panellbl] at (3.0,0.75) {Direct autograd: every $u_i$ stored, $O(N_t)$ memory};
  \begin{scope}[yshift=-2.6cm]
    \node[adjnd] (a0) at (0,0) {$\mathbf{a}(t_0)$};
    \node[adjnd] (a1) at (6.0,0) {$\mathbf{a}(t_1)$};
    \draw[arr, ink!65] (a1) -- (a0);
    \node[font=\scriptsize\sffamily\color{ink!70}] at (3.0,0.3) {adjoint ODE (Eq.~\ref{eq:adjoint}) integrated backward};
    \node[font=\scriptsize\sffamily\color{ink!70}] at (3.0,-0.3) {$u(t)$ recomputed as needed, nothing stored};
    \node[panellbl] at (3.0,0.75) {Adjoint sensitivity: $O(1)$ memory, independent of $N_t$};
  \end{scope}
\end{tikzpicture}}
\caption{Direct autograd (top) keeps every intermediate state $u_i$ in memory
so the backward pass can reuse it, giving $O(N_t)$ memory. Adjoint
sensitivity (bottom) instead solves the backward-time equation of
Equation~\ref{eq:adjoint}, recomputing $u(t)$ as needed, giving $O(1)$
memory regardless of $N_t$.}
\label{fig:adjoint-schematic}
\end{figure}

\subsection{A learnable hybrid basis}
\label{sec:hybrid-method}

Rather than choosing one basis function $B_i(x)$ in
Equation~\ref{eq:kan-edge} in advance, we can make the choice itself a
trained quantity: a learnable convex blend of two candidate bases inside
every edge,
\begin{equation}
\phi(x) = \alpha\cdot\mathrm{Spline}(x) + \beta\cdot\mathrm{RBF}(x),
\qquad (\alpha,\beta) = \operatorname{softmax}(\text{logits}),
\label{eq:hybrid-basis}
\end{equation}
trained end-to-end from a neutral start ($\alpha=\beta=0.5$) with no
hand-set schedule (Figure~\ref{fig:hybrid-schematic}). We produce the blend
weights with a softmax over two learnable logits rather than, say, a single
sigmoid-gated weight, because a softmax guarantees $\alpha+\beta=1$ by
construction. The two bases then always form a genuine convex combination,
so the gate can only trade one basis off against the other instead of
scaling the whole edge output up or down as a side effect. One pair of
logits is shared across every edge, so the model learns a single global
blend. This keeps the number of new parameters at two; it is a
simplification, not a claim that per-edge blending would fail.

\begin{figure}[t]
\centering
\resizebox{\columnwidth}{!}{%
\begin{tikzpicture}[node distance=10mm, font=\small,
  boxnd/.style={draw=ink, fill=accentsoft, minimum size=7mm, rounded corners=2pt},
  fnbox/.style={draw=ink!60, fill=codebg, minimum width=20mm, minimum height=7mm, rounded corners=2pt},
  gatend/.style={draw=accent, fill=white, circle, minimum size=6mm, inner sep=0pt, font=\bfseries\small\color{accent}},
  sumnd/.style={draw=accent, fill=accentsoft, circle, minimum size=7mm, inner sep=0pt, font=\bfseries\color{accent}},
  arr/.style={-{Stealth[length=2mm]}, draw=ink!65, thick}
]
  \node[boxnd] (x) at (0,0) {$x$};
  \node[fnbox] (sp) at (2.6,1.1) {$\mathrm{Spline}(x)$};
  \node[fnbox] (rbf) at (2.6,-1.1) {$\mathrm{RBF}(x)$};
  \node[gatend] (gA) at (5.0,1.1) {$\times$};
  \node[gatend] (gB) at (5.0,-1.1) {$\times$};
  \node[sumnd] (sum) at (7.2,0) {$+$};
  \node[boxnd] (phi) at (9.2,0) {$\phi(x)$};
  \node[fnbox, minimum width=19mm, minimum height=6mm, font=\scriptsize] (softmax) at (2.6,0) {$\operatorname{softmax}(\text{logits})$};
  \draw[arr] (x) -- (sp);
  \draw[arr] (x) -- (rbf);
  \draw[arr] (sp) -- (gA);
  \draw[arr] (rbf) -- (gB);
  \draw[arr] (gA) -- (sum);
  \draw[arr] (gB) -- (sum);
  \draw[arr] (sum) -- (phi);
  \draw[arr, accent!85] (softmax.north) -- node[right, font=\scriptsize\sffamily\color{accent!85}, pos=0.7]{$\alpha$} (gA.south);
  \draw[arr, accent!85] (softmax.south) -- node[right, font=\scriptsize\sffamily\color{accent!85}, pos=0.7]{$\beta$} (gB.north);
\end{tikzpicture}}
\caption{The hybrid-basis edge of Equation~\ref{eq:hybrid-basis}: every edge
evaluates both bases in parallel, and a shared softmax gate produces the
scalar weights $\alpha,\beta$ that multiply each basis's output before they
are summed into $\phi(x)$.}
\label{fig:hybrid-schematic}
\end{figure}

\subsection{Sparse regression and edge pruning}
\label{sec:sparse-method}

As a baseline from a completely different modelling approach, we use Sparse
Identification of Nonlinear Dynamics (SINDy) \cite{brunton2016sindy},
through the PySINDy library \cite{kaptanoglu2022pysindy}. SINDy fits a
governing equation as a sparse linear combination of candidate terms from a
fixed library; ours is every polynomial in $x$ and $y$ up to degree two,
$\{1,x,y,x^2,xy,y^2\}$. The coefficients are estimated by sequentially
thresholded least squares: fit by ordinary least squares against
numerically estimated derivatives, zero out every coefficient below a fixed
threshold, and refit on the surviving terms until the active set stops
changing. This is what makes the recovered model sparse and, in principle,
readable: unlike the KAN's 240 free parameters, SINDy returns a short list
of surviving terms with numeric coefficients that can be compared directly
with the true equation's. We hold the threshold fixed across the whole noise
sweep so that the comparison measures noise sensitivity rather than
per-level tuning.

The closest KAN-ODE analogue is to encourage sparsity during training and
prune weak edges afterwards. The base paper does this by adding an $L_1$
penalty on the spline coefficients to the loss,
\begin{equation}
\mathcal{L} = \mathrm{MSE}(\mathbf{u}_{\text{true}}, \mathbf{u}_{\text{pred}}) + \lambda\sum_{i} |c_i|,
\label{eq:l1-sparsity}
\end{equation}
where $c_i$ ranges over every spline coefficient in
Equation~\ref{eq:kan-edge} and $\lambda$ trades trajectory accuracy against
how many edges the penalty drives toward zero. Every checkpoint we trained
used $\lambda=0$, the plain mean-squared-error term alone. To prune, we rank
each edge by its mean spline-coefficient magnitude and mask the weakest
(Section~\ref{sec:track-e}).
